\documentclass[8pt]{strategynet-journal}
\usepackage{volt-paper-style}
\hypersetup{pdftitle={Pricing and Hedging a Multi-Asset Option with Volt},pdfauthor={research@strategynet.ai}}
\newcommand{\E}{\mathbb E}
\newcommand{\Q}{\mathbb Q}
\SNHeader{strategynet.ai:Multi-asset option learning}
\title{Pricing and Hedging a Multi-Asset Option with Volt}
\author{\href{mailto:research@strategynet.ai}{research@strategynet.ai}}
\date{October 2026}
\begin{document}
\maketitle
\begin{abstract}
\normalsize An option linked to several stocks can become most exposed to whichever stock leads the group. A hedge must respond as that exposure moves between names, and early exercise can bring the contract to an end before expiry. Valuing the remaining rights and financing the hedge therefore require a joint model of the underlying prices and a consistent account of exercise and trading. We develop this application in Volt using a Bermudan call on the largest of five stock prices. Neural continuation functions determine the holder's exercise decisions; independent simulations then estimate its payment value and a policy-tail upper valuation. The selected holder can also supply settlements to a writer that learns vector hedge holdings through the native application interface. A declared reconstruction of a published pricing experiment provides the numerical comparison. Three independently initialized models recover prices between 26.1576 and 26.1609 against the published 26.154, with wider confidence intervals. The application exposes the market model, backward learning calculation, valuation estimator and financed account through runnable interfaces, with independent numerical controls for CPU, Metal and CUDA execution.
\end{abstract}
\SNtwocolumntableofcontents
\begin{SNcolumns}

\section{The contract and market model}
Consider a holder who receives the amount by which the best-performing stock
exceeds a strike. With $d$ underlying prices $S^1,\ldots,S^d$ and strike $K$,
the cash payment on exercise is
\[
 g(S_t)=\left(\max_{1\leq i\leq d} S_t^i-K\right)^+.
\]
A different stock can lead at each date. That changes the distribution of
future payments and the stock positions useful to the writer. The Bermudan
contract adds a decision for the holder: exercise now or retain the remaining
opportunities. Our application follows this decision through valuation and
then into the writer's trading account.

The comparison uses Becker, Cheridito and Jentzen's five-asset case~\cite{american}.
All stocks start at 100, strike is 100, maturity is three years, and nine equal
exercise intervals give a four-month calendar. Each stock has annual volatility
20\% and continuous dividend yield 10\%; the interest rate is 5\%. Under the
pricing law $\Q$,
\[
 dS_t^i=(r-q_i)S_t^i\,dt+\sigma_i S_t^i\,dW_t^i.
\]
Thus the ex-dividend drift in this case is $-5\%$. The published row uses
independent Brownian components. The same generator accepts a supplied
correlation matrix and heterogeneous stock parameters.

For a time increment $\Delta t$, the exact GBM transition is
\[
 S_{t+\Delta t}^i=S_t^i\exp\!\left[
 (r-q_i-\tfrac12\sigma_i^2)\Delta t+
 \sigma_i\sqrt{\Delta t}\,Z_i\right].
\]
The vector $Z$ is standard Gaussian with the specified correlation. A
semidefinite Cholesky factor permits perfectly correlated assets as well as
full-rank inputs. Checked reconstruction and roundoff diagnostics accompany
the factorization. Indexed random streams bind every stock path to its seed,
namespace, global path ID and dates, so streaming batches retain their identities.

A portfolio application supplies different liabilities through the same
financial interface. The max-call is useful here because the published
benchmark supplies a concrete pricing target, while the underlying joint
process and vector account exercise the multi-asset implementation.

\section{Backward continuation learning}
Write $G_n=e^{-rt_n}g(S_{t_n})$ for a payment discounted to time zero. Suppose
a later exercise rule $\tau_{n+1}$ has already been chosen. The continuation
function at date $n$ estimates
\[
 c_n(x)=\E^{\Q}[G_{\tau_{n+1}}\mid S_{t_n}=x].
\]
The regression network minimizes squared error against the future rule's
realized discounted payments. Once fitted, it exercises where
$G_n\geq c_n(S_{t_n})$. Moving backward from expiry supplies the later rule
needed to fit each earlier date. All decisions use the observed state when
applied to a new path.

The state contains the five raw stock prices and the current discounted unit
payoff. Two hidden tanh layers each have 55 units; a linear output gives the
continuation value in currency per option unit. Batch normalization acts on
the input and hidden affine outputs. The final fitted date receives Xavier
initialization. Each earlier date starts from the next date's weights and
normalization statistics, with fresh Adam moments.

The article reports a declared reconstruction of the published settings.
The paper specifies the architecture, batch size, total update counts and
learning-rate values. We record the remaining choices explicitly: normalization
placement and population moments, epsilon 0.001, running-statistic weight
0.01 for new observations, Xavier-uniform kernels, finite path reuse, and
Adam's bias-corrected denominator with epsilon $10^{-8}$.

The final date uses 6,000 updates, with rates 0.1, 0.01, 0.001 and 0.0001 for
1,000, 2,000, 2,000 and 1,000 updates. Earlier dates use 3,500 updates with
500, 1,000, 1,500 and 500 updates at those rates. Batch size is 8,192.
Each model uses 1,048,576 training paths and 65,536 separate diagnostic
validation paths. Row sampling is uniform with replacement from resident
training tensors.

The scheduled final snapshot at every date supplies the primary exercise rule
and earlier training targets. Separate validation snapshots remain diagnostic.
This fixes selection before fresh pricing paths are generated. Completed-date
archives retain the market identities and selected future networks, allowing
an interrupted fit to restart the next earlier date reproducibly.

\begin{center}\begin{minipage}{\linewidth}\centering
\begin{tikzpicture}[font=\normalsize,>=Latex,x=1cm,y=1cm]
\draw[->] (0,0)--(6.4,0);
\foreach \x/\lab in {.4/{$0$},1.8/{$t_1$},3.3/{$t_2$},5.9/{$T$}}
 {\fill (\x,0) circle(1pt);\node[below] at(\x,-.08){\lab};}
\node[above,align=center] at(5.9,.35){Known\\payoff};
\draw[->,darkblue,thick] (5.5,.65)--(3.8,.65);
\draw[->,darkblue,thick] (3,.65)--(2.2,.65);
\node[above,align=center,darkblue] at(3.1,1.0){Fit continuation and exercise\\from later selected decisions};
\node[below,align=center] at(3.2,-.7){Apply the selected policy forward on fresh paths};
\end{tikzpicture}
\captionof{figure}{Training proceeds backward. Pricing and trading follow
new paths forward, preserving the information available at each decision.}
\end{minipage}\end{center}

\section{Independent policy-tail valuation}
A selected holder produces an expected discounted payment
$L=\E^{\Q}G_{\tau}$. Its admissible stopping rule makes $L$ a lower value
relative to optimal exercise. A sample mean estimates $L$ with Monte Carlo
uncertainty. An upper expectation can be constructed from the policy's value
martingale using nested conditional simulation~\cite{stopping,american}.

At each outer state $X_n$, simulate independent future paths that continue
past the current date and apply the selected holder from $n+1$. Their average
payment $\overline C_n$ estimates the continuation of that actual future rule.
If $f_n$ is the binary exercise decision on the outer state, define
\[
 \Delta\widetilde M_n=f_nG_n+(1-f_n)\overline C_n-\overline C_{n-1},
 \qquad \overline C_N=0.
\]
Starting from $\widetilde M_0=0$, accumulate these increments from date one.
Each conditional mean is reused in its two adjacent increments. The upper
estimate averages $\max_n(G_n-\widetilde M_n)$ along the full outer path,
including exercise opportunities after the holder's first stopping date.

This noisy process is the policy-value martingale plus conditionally centered
simulation error relative to the outer observations. The paper's conditional
Jensen argument supplies the upper expectation. The conditional simulations
retain their own random addresses for each parent state, date and child;
each outer observation also has its own initial continuation estimate.
Pooling several parent groups into a device batch preserves those addresses
and each parent's payoff summation.

The assessment uses 4,096,000 lower-estimate paths, 2,048 upper outer paths and
2,048 conditional branches per state. Three selected models share common
assessment market streams, allowing their realized payments to be compared.
Training, validation, outer and conditional sampling have separate declared
random domains. Reported intervals use the outer-sample standard errors and
an asymptotic joint 95\% lower/upper enclosure. The lognormal max-call is
unbounded, so this calculation uses its finite-variance sampling assumptions.

For the chosen source row, the published lower and upper estimates are 26.156
and 26.152, with midpoint 26.154 and enclosure $[26.138,26.171]$. Finite
sample estimates can cross even though their population expectations obey
the lower/upper inequalities. Our reports retain their observed ordering.
Before fitting, we set two numerical comparison targets: absolute midpoint
difference at most 0.05 and enclosure width at most 0.10. Both proximity and
precision enter the assessment.

\section{The writer's vector account and application API}
The selected holder can supply settlements to a writer. For continuous yield
$q_i$, use discounted dividend-reinvested security values
\[
 P_n^i=e^{(q_i-r)t_n}S_{t_n}^i.
\]
These share the time-zero currency basis of $G_n$. Cash growth is one, and a
holding $h_n^i$ corresponds to $h_n^i e^{q_i t_n}$ physical shares. Reinvestment
is included in the security definition.

Let $w_n$ be wealth before rebalancing and let $\lambda_i$ be proportional
stock costs. Choosing target holdings $h_n$ gives arrival wealth
\[
 w_{n+1}^{-}=w_n-
 \sum_i\lambda_iP_n^i|h_n^i-h_{n-1}^i|
 +\sum_i h_n^i(P_{n+1}^i-P_n^i).
\]
At exercise, the account liquidates its hedge, charges the declared closing
cost and pays the signed discounted liability once. Later holdings remain
zero. Quantity and multiplier scale the payment once; the caller supplies
initial capital for the complete position.

The native writer observes current stocks, discounted unit payoff and previous
holdings, then chooses a vector of continuous targets. Its complete-account
squared-loss objective differentiates through trades, fees and settlement.
The supplied writer example uses date-specific ReLU networks and this terminal
account objective. The paper's interval hedge method instead learns changes
in approximate continuation value using reinvested-stock features.

\texttt{maxcall\_episodes} verifies the actual holder checkpoint's digest,
applies its first admissible exercise decision and prepares native scenario
batches. The existing training, evaluation, holdings-inference and checkpoint
operations accept those batches. Jormungandr can supervise the same Volt
worker through its native HTTP service~\cite{api}.
\end{SNcolumns}
\noindent\begin{minipage}{\textwidth}
\begin{lstlisting}[basicstyle=\ttfamily\normalsize,frame=single,breaklines=true]
cargo build --release -p volt-compute --bin maxcall_learning --bin maxcall_episodes --bin native_learner
python3 examples/maxcall_application.py --output /tmp/volt-maxcall-example
\end{lstlisting}
\end{minipage}
\begin{SNcolumns}
The example fits a holder, prepares independent writer training/validation
ranges, selects a writer checkpoint and assesses its restored version on
fresh paths. Inputs, costs, normalization, market identities and model
snapshots accompany the outputs. A small default run demonstrates the full
sequence, while the published-case study uses the larger declared budgets.
The background primer develops the contract and account arithmetic~\cite{primer}.

\section{Precision and device performance}
The primary price comparison uses CPU float64. Independent calculations check normalization,
all parameter gradients, optimizer moments and post-update inference.

Single-precision tests identified a sensitive normalization representation.
A constant hidden affine bias is removed by the following batch normalization,
and an input-normalization translation is likewise removed by the next
normalized affine layer. In exact arithmetic these are redundant training
directions. Small floating-point cancellation residuals can be amplified by
Adam, while stored running means adjust over time.

Volt exposes an optional fixed-zero representation for these translations.
Converting an existing network also adjusts the following running means, so
its represented inference function is preserved. The remaining biases and
normalization shifts retain their learned roles. Independent Keras and NumPy
controls check both the converted and already-fixed forms on CPU, Metal and
CUDA. The original raw comparison and the stabilized comparison remain
separate in the evidence records.

The performance workload uses the same network, inputs, normalization,
arithmetic and precision across its compared devices. Each block restores
common weights and statistics. Warm-up updates precede synchronized measured
updates; simulation, upload, first-use setup and readback retain separate
measurements. CPU runtimes use one thread. Device memory requirements are
checked before allocation and recorded alongside actual runtime telemetry.
The implementation builds on Burn's tensor/autodiff interfaces and CubeCL
execution, using the repository's pinned versions~\cite{burn}.

\section{Pricing results and execution times}
All three independently initialized models passed the comparison targets.
Their midpoint estimates were 26.1576, 26.1609 and 26.1589, respectively,
against the published 26.154. The absolute differences ranged from 0.0036 to
0.0069, and each joint enclosure had width about 0.0745. The source row's
published enclosure has width 0.033; the present estimates therefore achieve
close price agreement with wider sampling uncertainty.

\begin{center}
\begin{tabular}{lrrr}
\toprule
Model & Lower & Upper & Midpoint\\
\midrule
Published &26.1560&26.1520&26.1540\\
Volt A &26.1393&26.1759&26.1576\\
Volt B &26.1427&26.1790&26.1609\\
Volt C &26.1407&26.1771&26.1589\\
\bottomrule
\end{tabular}
\end{center}
The respective joint 95\% enclosures were $[26.1206,26.1950]$,
$[26.1239,26.1984]$ and $[26.1221,26.1964]$. Common assessment paths couple
the three results; the three initializations measure training variability
under that shared assessment. Every checkpoint was selected and archived
before the first assessment result was opened.

Each full CPU float64 fit completed 30,500 updates in about 770 seconds.
The three fits ran concurrently, with one CPU thread per fit. This duration
includes the complete backward training sequence and its progress checks;
independent pricing follows as a separate calculation. The terminal-only
European value is 23.0516, obtained by independent integration of the
maximum-price distribution. Early exercise adds about 3.11 in this case.

The optional fixed-translation representation passed twelve independent
controls spanning CPU float64/float32, Metal float32 and CUDA float32, with
both converted and already-fixed initial states. The original float32
representation exposed post-update discrepancies in the redundant
translations; those observations motivated the explicit alternative.


The matched float32 measurements below use the optional fixed-translation
representation and strict native-dtype products. All eighteen CPU/GPU workload
comparisons passed the numerical controls. Each entry is the median of sixty
measured updates across three restored blocks, each preceded by three warm-ups.
Update times include forward/backward calculation, Adam, synchronization and
scalar checks. Here ms means milliseconds.

\begin{center}
\setlength{\tabcolsep}{3pt}
\begin{tabular}{rrrrrr}
\toprule
& &\multicolumn{2}{c}{Apple M3 Max}&\multicolumn{2}{c}{Linux workstation}\\
$d$&Batch&CPU&Metal&CPU&CUDA\\
\midrule
5&256&0.65&7.08&1.39&5.38\\
5&1,024&2.33&9.67&4.34&3.05\\
5&8,192&16.62&38.02&38.48&11.91\\
10&256&0.73&6.79&1.50&2.89\\
10&1,024&2.44&8.79&5.45&3.14\\
10&8,192&18.57&41.94&41.97&11.19\\
50&256&1.34&8.66&2.38&2.76\\
50&1,024&4.54&10.29&9.48&3.56\\
50&8,192&34.89&56.68&77.69&16.13\\
\bottomrule
\end{tabular}
\end{center}
The Mac runs macOS 27.0.1 with 128 GiB unified memory. The Linux system uses
one CPU thread on a Threadripper PRO 7985WX and an RTX 4090 with 24 GiB device
memory, CUDA 12.5 and driver 555.42.02. Both use Rust 1.99.1 and Burn 0.22.1.
Metal fusion is enabled; CUDA uses the qualified unfused configuration.
Background host utilization was recorded throughout, including concurrent
CPU training on the Mac. These are measurements under that observed load.

CPU execution was faster than Metal throughout these training cases. CUDA
became useful at batch 1,024: the five-stock case reached a 1.42-fold update
speed ratio over its same-host CPU, increasing to 3.23-fold at batch 8,192.
The fifty-stock, batch-8,192 ratio was 4.81. A larger state and resident batch
provide enough arithmetic to amortize device dispatch and synchronization.

First-use work also matters. For five stocks and batch 8,192, the first training
update took about 629 ms on Metal and 2,039 ms on CUDA. First inference and
readback took 103 and 412 ms respectively, while warmed inference/readback took
1.62 and 0.63 ms. These first-touch figures include the existing compiler/cache
state. The full timing record separates parameter setup, upload, first use and
steady updates, making it possible to choose a device for the actual request.


Independent tests cover market moments, financed cash flows, conditional
sampling, completed-date recovery and executable checkpoint provenance.

This study establishes one five-asset pricing comparison and a runnable vector
hedging application. Other contract parameters and hedge objectives can be
assessed through the same interfaces. Factor forecasts and live vendor inputs
require their own calibration and observation-timing studies.
\end{SNcolumns}
\clearpage
\begin{SNcolumns}
\begin{thebibliography}{99}
\bibitem{american} S. Becker, P. Cheridito and A. Jentzen. Pricing and hedging American-style options with deep learning. \emph{JRFM} 13(7):158, 2020. \url{https://arxiv.org/abs/1912.11060v3}.
\bibitem{stopping} S. Becker, P. Cheridito and A. Jentzen. Deep Optimal Stopping. \emph{JMLR} 20(74):1--25, 2019. \url{https://jmlr.org/papers/v20/18-232.html}.
\bibitem{api} Volt. Multi-asset market, learning and complete writer application guides. \url{https://github.com/ingvor-restrat/volt/blob/main/docs/multi-asset-maxcall.md}. \url{https://github.com/ingvor-restrat/volt/blob/main/docs/maxcall-learning.md}. \url{https://github.com/ingvor-restrat/volt/blob/main/docs/maxcall-application.md}.
\bibitem{primer} StrategyNet. \emph{Options, Exercise and Learned Hedging: A Practical Primer}. 2026. \url{https://github.com/ingvor-restrat/volt/blob/main/docs/latex/compute/volt_option_learning_primer.pdf}.
\bibitem{burn} Burn project. Tensor, automatic-differentiation and execution architecture. \url{https://burn.dev/books/burn/}.
\end{thebibliography}
\end{SNcolumns}
\end{document}
